% Section 11.2, from lectures/L11/appendix/b_exercises.md.

\section{Uppgifter}
\label{c11:sec:uppgifter}

\begin{aside}
\textbf{Så arbetar du med uppgifterna.} Uppgifterna repeterar kapitel~\ref{c7:ch}--\ref{c10:ch}
inför övningsduggan och görs med fördel på egen hand, en del per kapitel. Börja med de delar du
känner dig minst säker på.

\facitnotis{L11}
\end{aside}

% ------------------------------------------------------------------------------------------
\uppgiftsdel{Del 1}{Vektorer}
Uppgifterna repeterar \kapref{7}.

\uppgift{1.1}{Räkning med vektorer}
Vektorerna $\mathbf{u} = (2;\,5)$ och $\mathbf{v} = (-3;\,1)$ är givna. Bestäm:
\begin{delar}(4)
  \task $\mathbf{u} + \mathbf{v}$
  \task $\mathbf{u} - \mathbf{v}$
  \task $3\mathbf{u}$
  \task $2\mathbf{u} - 3\mathbf{v}$
\end{delar}

\uppgift{1.2}{Belopp och vinkel}
\begin{parts}
  \item Beräkna $\abs{(8;\,6)}$.
  \item Beräkna $\abs{(-5;\,-12)}$.
  \item Bestäm vinkeln mot positiv $x$-axel för $(8;\,6)$.
  \item Bestäm vinkeln för $(-5;\,-12)$. Tänk på kvadrantkorrigeringen.
\end{parts}

\uppgift{1.3}{Skalärprodukt}
\begin{parts}
  \item Beräkna $\mathbf{u} \cdot \mathbf{v}$ för $\mathbf{u} = (3;\,1)$ och
    $\mathbf{v} = (2;\,-6)$.
  \item Vad säger resultatet om vektorerna?
  \item Beräkna vinkeln mellan $\mathbf{u} = (1;\,2)$ och $\mathbf{v} = (3;\,1)$.
\end{parts}

\uppgift{1.4}{Enhetsvektor och given längd}
Vektorn $\mathbf{v} = (-8;\,6)$ är given.
\begin{parts}
  \item Bestäm enhetsvektorn i riktning $\mathbf{v}$.
  \item Bestäm en vektor med längden $25$ i samma riktning som $\mathbf{v}$.
  \item Bestäm en vektor med längden $5$ som är motriktad $\mathbf{v}$.
\end{parts}

\uppgift{1.5}{Spänningar som visare}
Över ett motstånd och en kondensator i serie ligger spänningarna
$\mathbf{U}_R = (9;\,0)\,\text{V}$ och $\mathbf{U}_C = (0;\,-12)\,\text{V}$.
\begin{parts}
  \item Beräkna den totala spänningen $\mathbf{U} = \mathbf{U}_R + \mathbf{U}_C$.
  \item Beräkna $\abs{\mathbf{U}}$.
  \item Beräkna vinkeln mot positiv $x$-axel.
  \item Varför är $\abs{\mathbf{U}}$ inte lika med $9 + 12 = 21\,\text{V}$?
\end{parts}

% ------------------------------------------------------------------------------------------
\uppgiftsdel{Del 2}{Funktioner}
Uppgifterna repeterar \kapref{8}.

\uppgift{2.1}{Definitions- och värdemängd}
Bestäm $D_f$ och $V_f$ för:
\begin{delar}(3)
  \task $f(x) = \sqrt{x - 2}$
  \task $f(x) = \dfrac{1}{x + 1}$
  \task $f(x) = -x^2 + 4$
\end{delar}

\uppgift{2.2}{Linjär modell för en termistor}
Resistansen hos en termistor beror linjärt på temperaturen enligt $R(T) = kT + m$. Mätning ger
$R(0^{\circ}\text{C}) = 200\,\Omega$ och $R(40^{\circ}\text{C}) = 120\,\Omega$.
\begin{parts}
  \item Bestäm $k$ och $m$.
  \item Beräkna $R$ vid $25^{\circ}\text{C}$.
  \item Vid vilken temperatur är $R = 100\,\Omega$?
  \item Är funktionen växande eller avtagande?
\end{parts}

\uppgift{2.3}{Exponentiell modell}
Antalet fel i ett system minskar med $15\,\%$ per månad. Vid start finns $80$ fel.
\begin{parts}
  \item Skriv en funktion $f(x)$ för antalet fel efter $x$ månader.
  \item Beräkna antalet fel efter $6$ månader.
  \item Efter hur många månader understiger antalet $20$ fel?
\end{parts}

\uppgift{2.4}{Effekt som potensfunktion}
Effekten i ett motstånd ges av $P(I) = 25I^2$ watt för $0 \leq I \leq 2\,\text{A}$.
\begin{parts}
  \item Beräkna $P(0{,}4)$.
  \item Beräkna $P(2)$.
  \item Bestäm $V_P$ för det angivna intervallet.
  \item Vid vilken ström är $P = 36\,\text{W}$?
\end{parts}

\uppgift{2.5}{Rationell funktion}
Funktionen $f(x) = \dfrac{x^2 - 9}{x - 3}$ är given.
\begin{parts}
  \item Faktorisera täljaren.
  \item Vilka $x$-värden är inte tillåtna?
  \item Förenkla uttrycket.
  \item Bestäm funktionens nollställe.
\end{parts}

% ------------------------------------------------------------------------------------------
\uppgiftsdel{Del 3}{Trigonometri}
Uppgifterna repeterar \kapref{9}.

\uppgift{3.1}{Vinkelmått}
\begin{delar}(2)
  \task Omvandla $135^{\circ}$ till radianer.
  \task Omvandla $-120^{\circ}$ till radianer.
  \task Omvandla $\dfrac{5\pi}{6}$ till grader.
  \task Omvandla $2\,\text{rad}$ till grader.
\end{delar}

\uppgift{3.2}{Egenskaper ur ekvationen}
En växelspänning ges av $u(t) = 6\sin\left(400\pi t - \dfrac{\pi}{3}\right)$ volt.
\begin{parts}
  \item Ange amplituden.
  \item Ange vinkelhastigheten och frekvensen.
  \item Ange periodtiden.
  \item Hur många millisekunder motsvarar fasen, och är kurvan tidig eller sen?
\end{parts}

\uppgift{3.3}{Skriv ekvationen}
En växelspänning har amplituden $12\,\text{V}$, periodtiden $T = 2\,\text{ms}$ och fasen
$30^{\circ}$. Bestäm $u(t)$ på formen $\abs{U}\sin(\omega t + \delta)$ med fasen i radianer.

\uppgift{3.4}{Momentanvärde, effektivvärde och effekt}
En växelspänning ges av $u(t) = 8\sin(100\pi t)$ volt.
\begin{parts}
  \item Beräkna $u(5\,\text{ms})$.
  \item Beräkna $U_{\text{RMS}}$.
  \item Spänningen läggs över ett motstånd $R = 50\,\Omega$. Beräkna effekten.
\end{parts}

\uppgift{3.5}{Trigonometrisk ekvation}
En växelspänning ges av $u(t) = 10\sin(200\pi t)$ volt.
\begin{parts}
  \item Vid vilken tidpunkt är $u(t) = 5\,\text{V}$ första gången?
  \item Ange nästa tidpunkt då $u(t) = 5\,\text{V}$.
  \item Vid vilken tidpunkt nås toppvärdet första gången?
\end{parts}

% ------------------------------------------------------------------------------------------
\uppgiftsdel{Del 4}{Logaritmer och decibel}
Uppgifterna repeterar \kapref{10}.

\uppgift{4.1}{Exponential- och logaritmekvationer}
Lös ekvationerna.
\begin{delar}(4)
  \task $2^x = 128$
  \task $10^x = 500$
  \task $e^{2x} = 5$
  \task $\log(x - 1) = 2$
\end{delar}

\uppgift{4.2}{Logaritmlagarna}
Beräkna eller förenkla.
\begin{delar}(4)
  \task $\log 4 + \log 25$
  \task $\log 500 - \log 5$
  \task $3\log 2$
  \task $\ln e^3$
\end{delar}

\uppgift{4.3}{Förstärkning i dB}
En förstärkare har $U_{\text{in}} = 20\,\text{mV}$ och $U_{\text{ut}} = 2\,\text{V}$.
\begin{parts}
  \item Beräkna förstärkningen i dB.
  \item Efter förstärkaren sitter ett filter som dämpar $12\,\text{dB}$. Beräkna den totala
    förstärkningen i~dB.
  \item Beräkna den totala linjära förstärkningen.
  \item Beräkna utspänningen efter filtret.
\end{parts}

\uppgift{4.4}{Effekt i dB och dBm}
Effektnivån i dBm anges relativt $1\,\text{mW}$:
\[
  P_{\text{dBm}} = 10\log_{10}\left(\frac{P}{1\,\text{mW}}\right)
\]
\begin{parts}
  \item En förstärkare har $P_{\text{in}} = 2\,\text{W}$ och $P_{\text{ut}} = 50\,\text{W}$.
    Beräkna förstärkningen i dB.
  \item En sändare har uteffekten $20\,\text{dBm}$. Ange effekten i mW.
  \item Hur många dBm motsvarar $1\,\text{W}$?
\end{parts}

\uppgift{4.5}{Urladdning av en kondensator}
Spänningen ges av $u(t) = 12e^{-t/(RC)}$ volt med $R = 22\,\text{k}\Omega$ och
$C = 100\,\mu\text{F}$.
\begin{parts}
  \item Beräkna tidskonstanten $\tau = RC$.
  \item Beräkna spänningen efter $3$ sekunder.
  \item Efter hur lång tid har spänningen halverats?
  \item Efter hur lång tid är spänningen nere i $1\,\text{V}$?
\end{parts}
